\chapter * {Úvod}

Měření počasí přímo v~místě, kde ho potřebujeme znát, je užitečné například při práci na zahradě nebo při rozhodování o~zalévání. Údaje z~veřejných předpovědí a vzdálenějších stanic totiž nemusí odpovídat tomu, co se děje u~konkrétního domu. Tato práce se zabývá návrhem, stavbou a naprogramováním chytré meteostanice, která tyto údaje měří přímo na místě. Skládá se z~venkovního modulu, který měří teplotu, vlhkost, vítr, srážky a intenzitu osvětlení, a z~vnitřního modulu, který tato data přijímá, doplňuje je o~vlastní měření teploty, vlhkosti a tlaku a zobrazuje je uživateli.

Práci zpracovávám jako maturitní práci a hlavním důvodem, proč jsem si vybral právě meteostanici, je možnost získat při jejím vypracování praktické zkušenosti. Projekt spojuje mechanickou konstrukci, elektroniku i programování, takže si na něm mohu vyzkoušet celý postup od návrhu až po funkční zařízení. Zároveň se dá rozdělit na samostatné části, jako je venkovní modul, vnitřní modul, komunikace mezi nimi a software, které lze postupně navrhovat, vyrábět a testovat. Na trhu je řada hotových meteostanic, které jsou kvalitně zpracované a nabízejí více, než dokážu vytvořit já. Nechci proto vytvořit lepší zařízení, ale ověřit si, že podobné řešení dokážu samostatně navrhnout, postavit a uvést do provozu.

Cílem práce je navrhnout, sestavit a naprogramovat funkční prototyp meteostanice, který měří teplotu, vlhkost, tlak, intenzitu osvětlení, rychlost a směr větru a úhrn srážek. Naměřené hodnoty a jejich historie se zobrazují na displeji a~na webové stránce a součástí zařízení je také jednoduché doporučení k~zavlažování. Funkce zařízení ověřuji v~reálných podmínkách. Největším přínosem práce jsou pro mě zkušenosti, které získám při řešení konkrétních problémů při návrhu, výrobě a oživení zařízení. Součástí práce je také technická dokumentace.